\phantomsection\label{fig:vol01:warring-qin:seven-states}
\section{空间总览：七国竞争与秦的东进通道}
\begin{center}
\begin{tikzpicture}[x=.83cm,y=.83cm,font=\small]
  \fill[paper] (0,0) rectangle (12,7);
  \draw[blue!55,line width=1pt] (.2,4.9) .. controls (3.4,4.6) and (6.4,5.2) .. (10.8,4.5);
  \node[text=blue!70!black,above] at (5.4,4.95) {黄河};

  \fill[dynasty!28] (1.0,2.4) -- (3.8,2.4) -- (4.2,4.8) -- (2.2,5.2) -- cycle;
  \node[align=center] at (2.5,3.65) {秦\\关中};
  \fill[green!20] (4.2,4.8) -- (6.3,4.6) -- (6.5,6.2) -- (4.8,6.4) -- cycle;
  \node at (5.4,5.55) {赵};
  \fill[orange!22] (6.5,4.5) -- (8.1,4.4) -- (8.3,5.9) -- (6.7,6.1) -- cycle;
  \node at (7.3,5.2) {燕};
  \fill[cyan!20] (8.2,3.7) -- (10.6,3.8) -- (10.7,5.5) -- (8.5,5.7) -- cycle;
  \node at (9.5,4.7) {齐};
  \fill[yellow!24] (4.0,3.4) -- (5.2,3.3) -- (5.4,4.5) -- (4.3,4.7) -- cycle;
  \node at (4.7,4.0) {魏};
  \fill[red!16] (5.1,2.7) -- (6.2,2.7) -- (6.3,3.8) -- (5.3,4.0) -- cycle;
  \node at (5.7,3.3) {韩};
  \fill[purple!18] (6.1,.7) -- (9.2,.8) -- (9.4,3.6) -- (7.0,3.8) -- cycle;
  \node at (7.7,2.1) {楚};

  \draw[-{Stealth[length=2mm]},ultra thick,color=dynasty] (3.4,3.8) -- (4.7,4.0);
  \draw[-{Stealth[length=2mm]},ultra thick,color=dynasty] (4.9,4.0) -- (5.8,3.35);
  \draw[-{Stealth[length=2mm]},ultra thick,color=dynasty] (5.9,3.3) -- (7.1,2.7);
  \node[below,text=dynasty] at (4.5,3.55) {秦东进};
\end{tikzpicture}

\smallskip
\small\textit{图  战国七国相对位置示意：突出关中秦国向中原推进的通道，以及魏韩赵、齐、楚等政治体的空间关系；疆界和路线均为约略表达。}
\end{center}
